\maketitle

\begin{abstract}
The systematic loss by firms of public tenders in which they participate for long periods and several times may indicate a cartel. This paper proposes a screening method to detect cartels in public tenders by considering frequent losers---firms identified as outliers through the interquartile range of losses among ``always loser'' bidders. Using administrative data on over 3.8 million transactions from the electronic procurement platform of S\~{a}o Paulo State, Brazil, from 2009 to 2019, we estimate that tenders with frequent losers are associated with approximately 10\% higher negotiated prices, 32\% more participating firms, and 29\% more bids. The simultaneous occurrence of higher prices and inflated competition indicators is inconsistent with competitive or tacit collusion models but consistent with explicit cartel behavior combined with a detection-avoidance strategy. These results are robust to alternative IQR thresholds, continuous treatment measures, multi-way clustering, coarsened exact matching, inverse probability weighting, and a staggered difference-in-differences design exploiting the timing of frequent losers' first entry into item markets. The proposed method addresses two limitations of traditional screening methods: (i) it distinguishes explicit from tacit collusion, and (ii) it enables ex-ante identification of suspicious behavior before the conclusion of public tender processes.
\end{abstract}

\medskip

\noindent \textbf{Keywords:} cartel; screening; public procurement; frequent losers; bid rigging.

\medskip

\noindent \textbf{JEL Codes:} D44, H57, K21, L41.

\bigskip

\noindent Darcio Genicolo-Martins\footnote{We are very thankful to Rita Joyanovic, Volnir Pontes Junior, M\'{a}rio Alexandre Reis da Silva, and the Department of Finance of S\~{a}o Paulo State (SEFAZ/SP) staff for their outstanding collaboration.}\textsuperscript{,}\footnote{Assistant Professor at Insper, Center for Regulation and Democracy. E-mail: \texttt{darciogm1@insper.edu.br}.}

\medskip

\noindent Paulo Furquim de Azevedo\footnote{Senior Fellow Research and Director of the Center for Regulation and Democracy at Insper. E-mail: \texttt{pfazevedo@insper.edu.br}.}
